\section{Introduction}

The traditional relational data model has long been the standard for data representation, but Big Data has revealed its limitations. A key challenge is the \textbf{V}ariety of Big Data, which encompasses multiple types and formats that originate from diverse sources and are inherently adherent to different models. There are structured, semi-structured, and unstructured formats; order-preserving and order-ignorant models; aggregate-ignorant and aggregate-oriented systems; etc. The naturally contradictory features of the so-called \textit{multi-model data}, combining distinct models with contradictory features, complicates all aspects of data management, from modelling to querying. 

Many tools now support multi-model capabilities -- e.g., more than two-thirds of the top 50 database management systems (DBMSs)\footnote{\url{https://db-engines.com/en/ranking}} align with Gartner's decade-old prediction of support for multiple models~\citep{Gartner2015}. However, the lack of standards for combining models leaves each DBMS with proprietary solutions. Even the variety of query languages~\citep{Jiaheng2023Survey} highlights the state of multi-model data systems.

Selecting the optimal tool for a specific use case is challenging, and extensive benchmarking is essential to compare tools across all target use cases. While single-model benchmarks and data generators exist for standard models, their adaptation to the multi-model paradigm is complex. Multi-model test cases must account for numerous combinations of different models using various strategies such as embedding, cross-model references, or redundancy. And this limits the availability and versatility of truly comprehensive multi-model benchmarks.~\citep{CSUR-Jiaheng,unibench}


Let us illustrate the problem with an example. We are building a web application that depends on various types of data. When it comes to selecting a database solution, we face a dilemma. We want the reliability of a relational database (e.g., for products and orders), but we also need to be able to perform complex graph queries (e.g., about the users, their mutual connections, and interactions with our products). We can leverage a multi-model database to store such data; however, we are uncertain whether this approach is more efficient than storing everything in a single database system. Many single-model datasets are similar to our expected data, so we can use them to benchmark the performance of the single-database solution. However, there are very few multi-model datasets available, nor there are tools for generating them.


To address this challenge, we introduce \textit{TransforMMer}, a tool for generating virtually any multi-model dataset. Instead of relying on traditional value generators with predefined distributions, our approach uses real-world single-model (or multi-model) datasets as input. It infers a schema of the input data, enables its user-specified modification to reflect the required combination of models, and transforms the output to the respective form. It also supports redundancy and versioning, which further increase its versatility. For its implementation, we integrate and extend selected tools from our toolset based on the unifying categorical representation of multi-model data~\citep{DBLP:journals/jbd/KoupilH22}. Combining and extending tools for modelling and transformations~\citep{mm_cat}, schema inference~\citep{mm_infer}, and evolution management~\citep{mm_evocat} with appropriate integration modules and interfaces allows us to ensure the transformation universally and correctly.

\paragraph{Outline} In Section~2, we review related work. Section~3 introduces the categorical representation of multi-model data. Section~4 introduces \textit{TransforMMer} and Section~5 its demonstration.

%%%%%%%%%%%%%%%%%%%%%%%%%%%%%%%%%%%%%%%%%%%%%%%%%%%%%

\section{Related Work}
\pLabel{sec:rel}
Two main approaches to benchmarking data management tools involve real-world data sets or synthetic data generators. Many focus on single models, but there are very few multi-model options.

Real-world datasets are primarily relational due to the dominance of relational DBMSs, accompanied by hierarchical models (e.g., JSON for NoSQL DBMSs) and graph data. Popular sources include Kaggle\footnote{\url{https://www.kaggle.com/}} or Harvard Dataverse\footnotei{\url{https://dataverse.harvard.edu/}}{.} Graph data can be found in collections like SNAP\footnotei{\url{https://snap.stanford.edu/data/}}{.} Governments (e.g., US\footnotei{\url{https://data.gov/}}{,} EU\footnotei{\url{https://data.europa.eu/}}{,} etc.) provide open data portals involving distinct data formats. Moreover, various datasets can also be found on GitHub or Google Dataset Search.

However, often we cannot easily find a suitable real-world dataset, and we need to use a data generator or benchmark. However, most are limited to specific data models, formats, or fixed use cases. For example, TPC-H and TPC-DS\footnote{\url{https://www.tpc.org/}} focus on the relational model. Benchmarks such as XMark~\citep{10.5555/1287369.1287455} are tailored to the document model. And also, many graph data generators exist~\citep{10.1145/3379445}. 

Regarding multi-model data, BigBench~\citep{10.1145/2463676.2463712} covers semi-structured and unstructured data, but lacks support for graph and array data models. UniBench~\citep{unibench} does not support the array data model either, and it considers only a single use case. M2Bench~\citep{10.14778/3574245.3574259} encompasses relational, document, graph, and array data models but also involves only three predefined use cases.

%%%%%%%%%%%%%%%%%%%%%%%%%%%%%%%%%%%%%%%%%%%%%%%%%%%%%%%%%%%%%

\section{Categorical Model of Data}
\pLabel{sec:cat}

First, to unify the terminology from different models, we use the following terms: A \textit{kind} corresponds to a class of items (e.g., a relational table or a collection of JSON documents). A \textit{record} corresponds to one item of a kind (e.g., a table row or a JSON document). A record consists of simple or complex \textit{properties} having their \textit{domains}.

We also recall the basic notions of category theory.
A \textit{category} $\mathbf{C} = (\mathcal{O}, \mathcal{M}, \circ)$ consists of a set of objects $\mathcal{O}$, set of morphisms $\mathcal{M}$, and a composition operation $\circ$ over the morphisms ensuring transitivity and associativity. Each morphism is modelled as an arrow $f : A \rightarrow B$, where $A, B \in \mathcal{O}$, $A = dom(f), B = cod(f)$. And there is an \textit{identity} morphism $1_A \in \mathcal{M}$ for each object $A$. A category can be visualized as a multigraph, with objects represented as vertices and morphisms as directed edges.

The core and the integration point of all our tools is an abstract representation of multi-model data called \textit{schema category}~\citep{DBLP:journals/jbd/KoupilH22}. It is defined as a tuple $\mathbf{S} = (\mathcal{O}_\mathbf{S}, \mathcal{M}_\mathbf{S}, \circ_\mathbf{S})$. \textit{Objects} in $\mathcal{O}_\mathbf{S}$ correspond to the domains of the ER model's entity types, attributes, and relationship types. 
Each schema object $o \in \mathcal{O}_\mathbf{S}$ is internally represented as a tuple $(key$, $label$, $superid$, $ids)$, where $key$ is an automatically assigned internal identity, $label$ is an optional user-defined name, $superid \ne \emptyset$ is a set of attributes (each corresponding to a signature of a morphism) forming the actual data contents a given object is expected to have, and $ids \subseteq \mathcal{P}(superid)$, $ids \ne \emptyset$ is a set of particular identifiers (each modeled as a set of attributes) allowing us to distinguish individual data instances uniquely.
%
\textit{Morphisms} connect appropriate pairs of objects. Each morphism $m \in \mathcal{M}_\mathbf{S}$ is represented as a tuple $(signature$, $dom$, $cod$, $label)$. The $signature$ allows us to distinguish all morphisms except the identity ones mutually. $dom$ and $cod$ represent the domain and codomain of the morphism. Finally, $label \in \{$ \texttt{\#property}, \texttt{\#role}, \texttt{\#isa}, \texttt{\#ident} $\}$ allows us to further distinguish morphisms with semantics \uv{has a property}, \uv{has an identifier}, \uv{has a role}, or \uv{is a}. The explicitly defined morphisms are denoted as \textit{base}, obtained via the composition $\circ$ as \textit{composite}.


The decomposition of a schema category $\mathbf{S}$, eventually partial or overlapping, is defined via a set of \textit{mappings}. For each kind, the mapping specifies where and how its records are stored in a selected single-/multi-model DBMS using a so-called \textit{access path}, which recursively describes the structure of a kind.

%Sample schema categories are provided in Figs.~\pRef{figure:sc-initial}--\pRef{figure:sc-multi}. 
For example, in Fig.~\pRef{figure:sc-initial}, we can see a sample schema category with kinds \textit{User}, \textit{Business} and \textit{Tip} having root objects denoted with green. They are mapped to JSON document collections from Fig.~\pRef{figure:data-initial}. For simplicity, we omit the most common edge label \uv{has a property} and replace uninteresting properties with dots in an oval.